% sec-03: closed-loop framing plus quantitative evidence.  Figure 2 is a
% mermaid-type state diagram of the loop with its human gate.
\section{A Closed Loop with a Human Gate}

The report calls closed-loop autonomous experimentation a strategic imperative
and notes that materials work taking around twenty years from laboratory to
deployment could be compressed by an order of magnitude
\cite{kratsios2026sciencegoldenage}. A clinical loop cannot be autonomous, and
this proposal does not ask for that. It asks for the loop to be closed
everywhere it can be, with one deliberate human gate that is measured rather
than assumed.

\begin{appfloat}[!tb]
\begin{appfig}[mermaid-type / state diagram]
% Five states around a loop, 3.1 apart horizontally and 2.3 vertically, with a
% guarded transition at the human gate.  Curved returns use looseness 0.8 so
% they cannot re-enter a state box.
\node[umlinit] (i) at (-1.9,0) {};
\node[mmsoft] (s1) at (0,0) {Observe\\intraoperative telemetry};
\node[mmsoft] (s2) at (3.5,0) {Propose\\advisory text};
\node[mmdec] (g) at (7.0,0) {Surgeon accepts?};
\node[mmgoal] (s3) at (10.4,0) {Execute\\approved motion};
\node[mmgray] (s4) at (5.25,-2.4) {Record decision,\\latency, and outcome};
\node[mmgray] (s5) at (10.4,-2.4) {Update the public dataset};
\draw[mmedgek] (i) -- (s1);
\draw[mmedgeb] (s1) -- (s2);
\draw[mmedgeb] (s2) -- (g);
\draw[mmedgeb] (g) -- node[mmlabel,pos=0.42] {yes} (s3);
\draw[mmedge] (g.south) -- node[mmlabel,pos=0.5] {no} (s4.north east);
\draw[mmedge] (s3) -- (s5);
\draw[mmedge] (s5) -- (s4);
\draw[mmedged] (s4.west) to[out=180,in=-90,looseness=0.8] (s1.south);
\node[font=\tiny\sffamily,text=protogray,anchor=north,text width=96mm,align=center]
  at (5.0,-3.4) {Both branches of the gate are recorded. A rejected
  recommendation is data, not a failure, and the rejection rate is a reported
  endpoint rather than an internal metric.};
\end{appfig}
\vspace{-0.7cm}
\figcaption{The loop, closed everywhere except at one measured human gate. Both
branches are recorded, so the rejection rate becomes a reported quantity rather
than an internal metric.}
\end{appfloat}

Quantitatively, the loop is already primed. An August 2025 quantitative systems
pharmacology run returned median overall survival of \textbf{12.8 months} for
the daraxonrasib combination against \textbf{5.4} for chemotherapy
\cite{qspmetpancre}; RASolute 302 reported \textbf{13.2} against \textbf{6.6
months} in the RAS G12 population nine months later \cite{rasolute302}. The
proximity is chronology, not validation: 1000 simulated against 241 enrolled, a
combination against a single agent, KRAS G12C against a primarily G12D and G12V
population. The digital-twin work carries an ASME V\&V 40 and ICH M15 aligned
credibility score of \textbf{81.9} over 55 verification notebook tests
\cite{daraxsims,asmevv40,ichm15}.
